
\section{1. Introduction}

\subsection{1.1 Motivation}

The softmax attention function (Vaswani et al., 2017) is the fundamental 
primitive of transformer architecture:

\[
\text{Attn}(Q, K, V) = \text{softmax}\!\left(\frac{Q K^\top}{\sqrt{d_h}}\right) V
\tag{1.1}
\]

Softmax has two properties that are convenient for optimization but awkward 
for interpretation and sparse computation:

1. \textbf{Dense output.} Every attention weight $A_{ij}$ is strictly positive. 
   Softmax never produces exact zeros. This makes attention patterns hard 
   to interpret -- a query spreads mass across all $L$ keys even when only 
   a few are relevant.

2. \textbf{No sparsity leverage.} Because every weight is nonzero, the attention 
   matrix must be materialized and multiplied densely. Sparse hardware 
   (e.g., NVIDIA 2:4 sparsity) cannot accelerate it.

\textbf{Empirical sparsity.} Prior work (Voita et al., 2019; Clark et al., 2019) 
shows that despite dense softmax outputs, the \textit{effective} attention is sparse: 
a small fraction of tokens receive most of the mass. \textbf{Figure 1.1} visualizes 
attention maps from a pretrained DistilBERT layer: the top 10\% of keys receive 
over 85\% of attention weight, while the bottom 50\% receive less than 2\%.

This discrepancy motivates \textbf{sparse attention activations}: replacing softmax 
with a function that produces \textbf{exact zeros} for low-scoring keys.

\subsection{1.2 The $\alpha$-Entmax Family}

Peters et al. (2019) introduce \textbf{$\alpha$-entmax}, a family of activation functions 
parameterized by $\alpha \in [1, 2]$:

\[
\text{$\alpha$-entmax}(s) = \arg\max_{p \in \Delta^{d-1}} \langle p, s \rangle + H_\alpha^T(p)
\tag{1.2}
\]

where $\Delta^{d-1}$ is the simplex and $H_\alpha^T$ is the Tsallis entropy 
with parameter $\alpha$. Two special cases are well-known:

\begin{itemize}
\item $\alpha = 1$: reduces to \textbf{softmax} (dense).
\end{itemize}
\begin{itemize}
\item $\alpha = 2$: reduces to \textbf{sparsemax} (Martins \& Astudillo, 2016),
\end{itemize}
  which projects scores onto the simplex and produces the sparsest 
  distribution.

Intermediate values interpolate between the two:

\begin{itemize}
\item $\alpha = 1.5$: \textbf{1.5-entmax}, which in practice produces ~40\% exact zeros
\end{itemize}
  while remaining differentiable (unlike sparsemax at boundaries).

This makes $\alpha = 1.5$ a natural default: sparse enough to matter, dense 
enough to train stably.

\subsection{1.3 Attention Sinks and Local Windows}

\textbf{Attention sinks.} Xiao et al. (2024) and others observe that the first 
few tokens in a sequence -- regardless of content -- receive disproportionate 
attention. This is an artifact of how transformers initialize positional 
information; preserving these tokens is critical for stability.

\textbf{Local windows.} Longformer (Beltagy et al., 2020) and BigBird (Zaheer et al., 
2020) restrict attention to a sliding window of radius $W$. This captures 
local syntax while keeping complexity $\mathcal{O}(L \cdot W)$ instead of 
$\mathcal{O}(L^2)$.

AdaSplash combines both: \textbf{local splash} attention keeps a window of radius 
$W$ around each query, plus the first 4 "sink" tokens, plus $\alpha$-entmax 
sparsity on top.

\subsection{1.4 Problem Statement}

Given a pretrained attention module with projections 
$W_Q, W_K, W_V, W_O \in \mathbb{R}^{d \times d}$, we seek a modified attention 
function $\mathcal{A}'$ such that:

1. \textbf{Sparse outputs.} $\mathcal{A}'$ produces \textbf{exact zeros} in attention 
   weights for low-scoring keys.
2. \textbf{Local structure.} Attention is restricted to a sliding window of radius 
   $W$ plus sink tokens.
3. \textbf{No base weight updates.} Base weights $W_Q, W_K, W_V, W_O$ are frozen.
4. \textbf{Trainable adapters.} LoRA rank-$r_L$ adapters are added for task adaptation.
5. \textbf{Multi-architecture support.} Works on flat (DistilBERT), nested 
   (RoBERTa/BERT), causal (GPT-2), GQA (Qwen), and RoPE.

\textbf{Note.} This is \textbf{not} a compression method. AdaSplash \textbf{increases} 
parameter count by the LoRA adapter size. It trades parameter overhead for 
attention sparsity and interpretability.

\subsection{1.5 Our Approach}

\textbf{AdaSplash} combines four components:

1. \textbf{$\alpha$-Entmax bisection.} We solve the entmax equation via bisection on 
   $\tau$, the Lagrange multiplier. The procedure runs for 25 iterations, 
   producing a numerically stable approximation with softmax fallback.

2. \textbf{Local splash masking.} For causal decoders, keys outside $|i - j| \leq W$ 
   are masked, along with the first 4 tokens preserved as attention sinks.

3. \textbf{Frozen base + LoRA.} Base projections are copied exactly from pretrained 
   weights. LoRA rank-$r_L$ adapters with $B = 0$ initialization provide 
   trainable capacity.

4. \textbf{Multi-architecture weight transfer.} A comprehensive \texttt{copy\\_weights} 
   handles GPT-2 Conv1D, RoBERTa/BERT nested, DistilBERT flat, and Qwen 
   GQA patterns.

\subsection{1.6 Contributions}

\begin{itemize}
\item \textbf{A stabilized $\alpha$-entmax implementation.} We provide a numerically robust
\end{itemize}
  bisection algorithm for $\alpha$-entmax with valid-token masking and softmax 
  fallback, available for $\alpha \in (1, 2]$. Section 4.2.
\begin{itemize}
\item \textbf{Local splash masking with sink preservation.} We combine sliding-window
\end{itemize}
  restriction with attention sink preservation (first 4 tokens). Section 4.3.
\begin{itemize}
\item \textbf{Architecture-agnostic injection.} The method works across flat, nested,
\end{itemize}
  causal, and GQA attention without modification. Section 5.
\begin{itemize}
\item \textbf{Systematic evaluation.} We benchmark on three models on SST-2 with
\end{itemize}
  ablations over $\alpha$, splash radius $W$, and LoRA rank. Section 7.
\begin{itemize}
\item \textbf{Honest disclosure.} We clearly state that AdaSplash \textbf{increases}
\end{itemize}
  parameter count due to LoRA, and report when the method is worth using. 
  Sections 8--9.
\begin{itemize}
\item \textbf{Full reproducibility.} Code, checkpoints, and evaluation scripts released
\end{itemize}
  under {{LICENSE}}.

\subsection{1.7 Paper Organization}

Section 2 reviews related work. Section 3 provides background on $\alpha$-entmax, 
attention sinks, and LoRA. Section 4 presents AdaSplash. Section 5 describes 
architecture integration. Section 6 details the experimental setup. Section 7 
reports results. Section 8 analyzes attention sparsity and ablation 
contributions. Section 9 discusses practical usage. Section 10 lists 
limitations. Section 11 concludes.



\subsection{Notation}
